\documentclass[11pt]{article}
\usepackage[margin=23mm]{geometry}
\usepackage[T1]{fontenc}
\usepackage{lmodern,amsmath,amssymb,amsthm,mathtools,booktabs,microtype,hyperref}
\hypersetup{colorlinks=true,linkcolor=blue,urlcolor=blue}
\newtheorem{theorem}{Theorem}
\newcommand{\Z}{\mathbb Z}
\allowdisplaybreaks
\setlength{\emergencystretch}{2em}
\title{CMF conservation and the arithmetic of\\$A\zeta(3)+B\zeta(4)$\\\large The trajectory $(16,17,-11,-12,-10;2,1,0,-1,-2)$}
\author{An arithmetic proof with exact rational certificates}
\date{10 October 2026}
\begin{document}\maketitle
\begin{abstract}
For the specified native MeijerG CMF, we establish a uniform common
divisor of the two prescribed zeta readouts. Taking their exponential
growth and approximation-error rates as given, the arithmetic estimate
implies $\liminf\delta_N>0.0113626$ for every fixed integer pair
$(A,B)\ne(0,0)$, and $\liminf\delta_N>0.2190118$ when $B=0$.
The proof uses an admissible reordered CMF path, a rational kernel with
double poles, an exact prime-carry partition, and a finite-field sum
that supplies additional common factors. These bounds do not assert
that the observed finite-depth deltas are the exact asymptotic values.
\end{abstract}

\section{The statement and the accepted analytic input}
Use $\operatorname{MeijerG}(5,2,5,5,1)$ with initial point
$(1,1,0,0,0;0,0,0,0,0)$ and the direction in the title. Put $n=N-1$,
exactly matching the user's indexing. The accepted initial-condition
matrix is
\[
 R_{A,B}=\begin{pmatrix}0&0&-6A&3A-2B\\12&0&0&0\end{pmatrix}.
\]
In all statements $A,B$ are fixed integers; they do not vary with $N$.
Let $r_N(A,B)$ be the quotient of the two entries of $R_{A,B}W_ne_0$.
Its reduced positive denominator is denoted $q_N(A,B)$.

The analytic input, used without reproving it, is
\[
 \log|Q_n|=\sigma n+o(n),\qquad
 \left|\zeta(s)-\frac{P_{s,n}}{Q_n}\right|
       \le \exp(-En+o(n)),\quad s=3,4,
\]
where the exact readouts $Q_n,P_{s,n}$ are identified below and
\[
 \sigma=53.287583322620205\ldots,\qquad
 E=112.807539247454680\ldots .
\]
Exact rational isolation and logarithm certificates for the associated
spectral constants accompany the proof. In particular
\[
 \sigma+30\log(3/2)<65.451536565866,\qquad E>112.807539247454.
\]
The full limiting matrix also has an eigenvalue $-1$; the selected
readouts' elimination of that mode is part of the given analytic input.
The triangle inequality gives the same error upper exponent $E$ for
every fixed $A\zeta(3)+B\zeta(4)$.

\begin{theorem}
Under this accepted analytic input, the reduced approximants satisfy
\[
 \boxed{\liminf_{N\to\infty}\delta_N(A,B)>0.0113626}
 \quad\text{for every fixed }(A,B)\in\Z^2\setminus\{(0,0)\}.
\]
If $B=0$, the stronger estimate is
\[
 \boxed{\liminf_{N\to\infty}\delta_N(A,0)>0.2190118}\qquad(A\ne0).
\]
Here $\delta_N=-1-\log|A\zeta(3)+B\zeta(4)-r_N|/\log q_N$ whenever
$q_N>1$ and the error is nonzero; an exact approximant may equivalently
be assigned $\delta_N=+\infty$. Fixed rational coefficients have the same
exponents after multiplication by one fixed denominator.
\end{theorem}

\section{What multiplication order contributes}
Conservation identifies products along admissible shifted lattice paths
with the same endpoints. It does not say that matrices evaluated at a
single vertex commute. We first decrease $a_2,a_3,a_4$, then increase
$a_0,a_1$, then increase $b_0,b_1$, and finally decrease $b_3,b_4$.
This order creates a long interval of triple zeros before the last
contour shift is needed. It also exposes factorial and harmonic blocks
that are invisible in an entrywise estimate of the original trajectory
matrix.

There is no requirement to use the same arithmetic certificate for the
two initial rows. Nor must one use the same path at different primes.
For a cut $w$, write $W_n=U_wV_w$. The identity
\[
 rW_ne_0=\sum_i(rU_w)_i(V_we_0)_i
\]
gives the valid certificate
\[
 v_p(rW_ne_0)\ge \min_i v_p\bigl((rU_w)_i(V_we_0)_i\bigr).
\]
One may maximize this independently for each $p$ and each initial row.
For the mixed numerator it is also useful to split
$r_{A,B}=Ar_3+Br_4$: bounds proved separately for $r_3$ and $r_4$ imply a
bound for every integer combination. This can be stronger than applying
one entrywise estimate to the already combined row.

The accompanying exact path experiments demonstrate this freedom in the
present CMF. They are not the justification for an asymptotic divisor.
The uniform argument below exposes one convenient rational readout and
then proves its primewise divisibility directly. It uses the conserving
family to explain the arithmetic, without claiming that an empirically
best ordering has been proved optimal.

% Uniform arithmetic readout for MeijerG(5,2,5,5,1).
% Here n=N-1; no asymptotic analytic assertion is used.
\section{A rational kernel for the two initial conditions}
Put $n=N-1\ge0$ and let $W_n$ be the native forward CMF product from
$(1,1,0,0,0;0,0,0,0,0)$ to
$(1+16n,1+17n,-11n,-12n,-10n;2n,n,0,-n,-2n)$.
Define
\[
 F_n(t)=\frac{\prod_{u=0}^2(t-(10+u)n)_{10n}}
 {(t+n)_{16n+1}(t+2n)_{14n+1}}.
\]
Its partial fractions are
\[
 F_n(t)=\sum_{j=n}^{17n}\left(\frac{a_j}{(t+j)^2}
                    +\frac{b_j}{t+j}\right),\qquad
 a_j=0\quad(j<2n\ \hbox{or}\ j>16n).
\]
The function is $O(t^{-2})$, so $\sum_jb_j=0$.
With $H_m^{(s)}=\sum_{k=1}^m k^{-s}$, set
\begin{align*}
 Q_n&=\sum_j a_j,\\
 P_{3,n}&=\sum_j\left(a_jH_{j+2n}^{(3)}+
                              \frac{b_j}{2}H_{j+2n}^{(2)}\right),\\
 P_{4,n}&=\sum_j\left(a_jH_{j+2n}^{(4)}+
                              \frac{b_j}{3}H_{j+2n}^{(3)}\right).
\end{align*}
Then the following identities hold uniformly:
\[
 Q_n=(-1)^n(W_n)_{00},\qquad (W_n)_{10}=0,
\]
\[
 P_{3,n}=(-1)^n\left(-\frac{(W_n)_{20}}2+
                              \frac{(W_n)_{30}}4\right),\qquad
 P_{4,n}=(-1)^n\left(-\frac{(W_n)_{30}}6\right).
\]
Consequently the requested approximant is exactly
$(A P_{3,n}+B P_{4,n})/Q_n$.

\paragraph{Proof by the CMF mechanism.}
For an integer vertex $w=(a_0,\ldots,a_4;b_0,\ldots,b_4)$ write
\[
 R_w(t)=(-1)^{a_0+a_1}
       \frac{\prod_{j=0}^4\Gamma(t+b_j)}
            {\prod_{i=0}^4\Gamma(t+a_i)}.
\]
On the path used below this is a rational function with poles of order at
most two, all at integers.  Initially it is $t^{-2}$, and at the endpoint
it is $(-1)^nF_n(t)$.  The gamma functional equation gives exactly the
native elementary CMF multipliers when $\theta$ is represented by $-t$.
To check the companion reduction, put
$H_w(t)=R_w(t)\prod_j(t+b_j)$.  Then
\[
 R_w(t)\prod_i(t+a_i-1)=H_w(t-1).
\]
Thus the differential equation reduces to a translate difference of
$H_w$, up to its immaterial overall sign.

Choose the coordinate-block order: first decrease $a_2,a_3,a_4$ to their
endpoint values; then increase $a_0,a_1$; then increase $b_0,b_1$; finally
decrease $b_3,b_4$.  During the first three blocks use a vertical contour
$c\in(0,1)$.  At their end the rational numerator has zeros of multiplicity
three at every integer $1,\ldots,10n$, so move the contour to
$c\in(2n,2n+1)$.  This move crosses no pole for either kernel
\[
 K_2(t)=\frac{\pi^2}{\sin^2\pi t},\qquad
 K_3(t)=\frac{\pi^3\cos\pi t}{\sin^3\pi t}.
\]
Use this latter contour for every remaining block.  Initially $H_w$ has
a zero of order at least three at $0$.  In the later blocks it has one at
$2n$: the possible loss of one zero when $b_4=-2n$ is restored by the
factor $t+b_4$ in $H_w$.  Therefore multiplication by either periodic
kernel and integration annihilates $H_w(t-1)-H_w(t)$.  Horizontal boundary
integrals vanish by the kernels' exponential decay.  This proves the
companion relation at every point used by the path; the elementary gamma
relations prove the edge relations.  Hence these integrals propagate by
the actual native CMF matrices, without a scalar normalization.

More precisely, the two initial rows of period integrals on the basis
$((-t)^r/t^2)_{r=0}^3$ are
\[
 (2\zeta(3),-\zeta(2),1,-1/2),\qquad
 (-3\zeta(4),\zeta(3),0,-1/2).
\]
The first two entries follow by closing the contour to the right.
For the last two, $K_2=-\pi(\cot\pi t)'$ gives
$\int K_2/(2\pi i)=1$; reflection about $1/2$ gives
$\int tK_2/(2\pi i)=1/2$; and $K_3=-K_2'/2$ gives the second row by
integration by parts.

There are also two purely rational functionals: the sum of the double-pole
coefficients and the sum of the simple-pole coefficients.  Both vanish on
integer translate differences and on polynomials.  Their initial rows are
$(1,0,0,0)$ and $(0,-1,0,0)$.  They therefore obey the same CMF transport.
At the endpoint they give $Q_n=(-1)^n(W_n)_{00}$ and $(W_n)_{10}=0$.
Finally, closing the endpoint contour to the right gives
\[
 \int K_2F_n\,\frac{dt}{2\pi i}
   =-\sum_{k=2n+1}^{\infty}F_n'(k)=2(Q_n\zeta(3)-P_{3,n}),
\]
\[
 \int K_3F_n\,\frac{dt}{2\pi i}
   =-\frac12\sum_{k=2n+1}^{\infty}F_n''(k)
   =-3(Q_n\zeta(4)-P_{4,n}).
\]
Comparing with the transported rows proves all asserted identities.
The case $n=0$ is immediate from $W_0=I$ and $F_0=t^{-2}$.

\section{A uniform denominator bound and the special $B=0$ saving}
Let $d_m=\operatorname{lcm}(1,\ldots,m)$ and, for $n\ge1$, put
\[
 C_n=\frac{(15n)!^2}{(10n)!^3}.
\]
Then
\[
 C_na_j\in\mathbb Z,\qquad d_{16n}C_nb_j\in\mathbb Z.
\]
In particular
\[
 6C_nd_{19n}^3(Q_n,P_{3,n})\in\mathbb Z^2,\qquad
 6C_nd_{19n}^4(Q_n,P_{3,n},P_{4,n})\in\mathbb Z^3.
\]
This removes one entire factor $d_{19n}$ when $B=0$, uniformly in $n$;
it is not an inference from finite gcd measurements.

\paragraph{Proof.}
Re-pair the denominator intervals as
$(t+n)_{15n+1}(t+2n)_{15n+1}$ and write
\[
 C_nF_n(t)=P(t)L_n(t)L_{2n}(t),\quad
 P(t)=\prod_{u=0}^2\binom{t-un-1}{10n},\quad
 L_c(t)=\frac{(15n)!}{\prod_{j=c}^{c+15n}(t+j)}.
\]
Every $L_c$ has integer simple-pole residues
$(-1)^{j-c}\binom{15n}{j-c}$.
The polynomial $P$ is integer-valued on integers.  Moreover
$d_{10n}P'(k)$ is integral at every integer $k$: apply the Newton-basis
identity
$\binom{x}{m}'=\sum_{i=1}^m(-1)^{i-1}\binom{x}{m-i}/i$
to each of its three binomial factors and use the product rule.
At a common pole the double coefficient is therefore integral, while its
simple coefficient is a sum of an integer multiple of $P'(-j)$ and
terms with integer numerator and denominator $j-k$.  At a pole of only one
factor only the latter kind occurs.  All such differences satisfy
$0<|j-k|\le16n$.  This proves both coefficient assertions, including the
simple poles outside the overlap.  Finally $j+2n\le19n$, and the harmonic
numbers of order $s$ are cleared by $d_{19n}^s$.

\section{Primewise contents of the whole readout}
We next divide the integer forms by a uniform product of primes. The
valuation of a sum is bounded after its full factorial and harmonic
structure has been exposed, rather than after each elementary matrix step.
Only primes
\[
 p>29,\qquad p^2>30n,\qquad p\nmid n
 \tag{E}
\]
are used. The omitted primes affect the logarithmic density by $o(n)$.
Put $x=n/p$ and $y=j/p$. Since every factorial argument is at most $29n$,
Legendre's formula has just its first term at these primes.
Define
\[
 T(x,y)=2\lfloor15x\rfloor-3\lfloor10x\rfloor+
 \sum_{u=0}^2\bigl(\lfloor y+(10+u)x\rfloor-\lfloor y+ux\rfloor\bigr).
\]
At a double pole, $2x\le y\le16x$, the exact valuation of $C_na_j$ is
\[
 s=T-\lfloor y-x\rfloor-\lfloor y-2x\rfloor
        -\lfloor16x-y\rfloor-\lfloor17x-y\rfloor.
 \tag{D}
\]
At the two simple-pole ranges it is the valuation of $C_nb_j$ that is
needed, and it is exactly
\begin{align*}
 b_-&=T+\lceil2x-y\rceil-1-\lfloor16x-y\rfloor
                   -\lfloor y-x\rfloor-\lfloor17x-y\rfloor,
           &&x\le y<2x,\\
 b_+&=T+\lceil y-16x\rceil-1-\lfloor y-2x\rfloor
                   -\lfloor y-x\rfloor-\lfloor17x-y\rfloor,
           &&16x<y\le17x.
\end{align*}
These formulas follow by evaluating the factorial products at each pole.
The ceilings retain the endpoint correction in factorials of the form
$(2n-j-1)!$; dropping this correction would invalidate some prime bounds.

For a double pole, $b_j=a_jL_j$, where $L_j$ is a signed sum of reciprocal
integers. Write $H_m=\sum_{k=1}^m1/k$ and set
\begin{align*}
 \ell(x,y)={}&\sum_{u=0}^2\left(H_{\lfloor y+ux\rfloor}
                           -H_{\lfloor y+(10+u)x\rfloor}\right)\\
 &+H_{\lfloor y-x\rfloor}+H_{\lfloor y-2x\rfloor}
  -H_{\lfloor16x-y\rfloor}-H_{\lfloor17x-y\rfloor}.
\end{align*}
Indeed, $pL_j$ reduces modulo $p$ to this rational number. All its
harmonic indices are smaller than $p$, by (E). Thus
\[
 v_p(C_nb_j)\ge s-\varepsilon(x,y),\qquad
 \varepsilon(x,y)=\begin{cases}0&\ell(x,y)=0,\\1&\ell(x,y)\ne0.\end{cases}
 \tag{LD}
\]
The equality test here is an equality in $\mathbb Q$, not a floating-point
or heuristic congruence test. If $\ell\ne0$ happens to vanish modulo a
particular prime, ignoring that cancellation is conservative.

Put $l=\mathbf1_{x>1/19}$ and $h=\mathbf1_{y+2x\ge1}$. The elementary
bound $v_p(H_{j+2n}^{(r)})\ge-rh$ gives the following lower bounds for
$6C_nd_{19n}^{r}P_{r,n}$, with $r=3,4$:
\[
 \begin{array}{c|c}
 \text{contribution}&\text{valuation lower bound}\\\hline
 \text{double coefficient times }H^{(r)}&s+rl-rh\\
 \text{simple coefficient at a double pole}&s-\varepsilon+rl-(r-1)h\\
 \text{left or right simple pole}&b_\pm+rl-(r-1)h.
 \end{array}
 \tag{V}
\]
The same minimum also bounds $6C_nd_{19n}^rQ_n$: its corresponding
valuation is $s+rl$, which is no smaller. Take the minimum over the three
specified $y$ ranges and then its maximum with zero, using the established
integer normalization for that last operation. Call the result $g_r(x)$.

\paragraph{An exact finite certificate for all depths.}
The attached \texttt{profiles/exact\_chambers.py} evaluates (D)--(V)
using rational arithmetic. Every breakpoint in $y$ is a line $y=m+kx$
with integer slope between $-12$ and $17$. Two such lines cross only at
a rational $x$ of reduced denominator at most $29$. The constant-floor
breakpoints and the harmonic cutoff are included in this list. Therefore
it suffices to check the $270$ open intervals of the Farey partition of
order $29$ in $(0,1)$, evaluating both the open $y$ intervals and their
endpoints. This is a complete finite partition of the formulas, not a
check at finitely many trajectory depths. Condition (E) prevents $x$ from
lying on its boundaries.

The output is
\[
 g_3(x)=g_4(x)=g(x)=\psi(\{x\})+h_0(x),
 \tag{G}
\]
where $\psi$ is the step function in the appendix and
\[
 h_0(x)=
 \begin{cases}
 3,&1/19<x<1/18,\\
 2,&1/18<x<1/17,\\
 1,&1/17<x<1/16,\\
 0,&\text{otherwise}.
 \end{cases}
\]
For $x\ge1$ the harmonic cutoff is always active. In (V) the double-pole
minimum is then $\min(s,s+1-\varepsilon)=s$, independent of $\ell$.
The other values are $b_\pm+1$. Each $y$ range has width at least one, and
these floor expressions are periodic under integral translations of $x$
and $y$. This proves the periodic tail in (G); no periodicity of the
harmonic expression $\ell$ is presumed.

It follows that the product of eligible primes $p^{g(n/p)}$ divides both
integer forms of orders three and four. Prime-dependent choices are
permitted because divisibility is certified separately in each local
ring $\mathbb Z_{(p)}$; unique factorization then combines the bounds.

\section{A uniform extra prime factor: $8n<p<11n$}
For every integer $n\ge1$ and prime $8n<p<11n$, the rational readout just
constructed satisfies
\[
 v_p(Q_n)\ge1,\qquad v_p(P_{3,n})\ge-1,\qquad
 v_p(P_{4,n})\ge-2.                                      \tag{BS}
\]
The termwise harmonic estimates give only $-2$ and $-3$ for the last two
valuations.  Thus (BS) supplies one additional common factor $p$ for both
initial conditions, uniformly throughout this entire prime interval.
This assertion does not require an asymptotic hypothesis or a finite-depth
extrapolation.

\paragraph{Root multiplicities.}
Write $F_n=\mathcal N/\mathcal D$ with the monic numerator and denominator
from its definition.  A numerator root lies in one of the three integer
intervals $[1,10n]$, $[n+1,11n]$, $[2n+1,12n]$.  A denominator root is the
negative of an integer in $[n,17n]$ or $[2n,16n]$.
For $0\le k<p$ let $c(k)$ and $d(k)$ count these roots congruent to $k$,
with multiplicity.  Direct interval counting gives
\[
 c(k)-d(k)\ge-1.
\]
Equality holds precisely when
\[
 \max(12n-p+1,\,2p-16n)\le k\le p-2n.                  \tag{S}
\]
On this support the three numerator roots are all $k$ itself, whereas the
four denominator roots are $k-p,k-p,k-2p,k-2p$.
In particular $k>2n$.

Here are details of the count, to make the chamber assertion uniform in
$n$ and $p$.  In the middle interval $2n<k\le p-2n$, all three roots $k$
are present, and both denominator roots $k-p$ are present.  There are at
most two additional denominator roots, both $k-2p$; they are both present
exactly when $k\ge2p-16n$.  No additional numerator root $k+p$ is present
exactly when $k>12n-p$.  This proves (S) on the middle interval.
For $p-2n<k\le p-n$ there are three roots on each side.  For $p-n<k<p$
there are two or three denominator roots and at least as many numerator
roots; a third denominator root can occur only when $p<10n$, when all
three numerator roots $k$ are present.
For $0\le k\le2n$ the denominator count is
$2+\mathbf1_{k\ge2p-17n}+\mathbf1_{k\ge2p-16n}$.
At $k=0$ the assertion follows immediately by comparing $p$ with
$17n/2$ and $10n$.  For $0<k\le n$, its third root forces $p\le9n$ and
hence a third numerator root, while its fourth forces $p\le17n/2$ and
hence a fourth numerator root.  For $n<k\le2n$, these same implications
are $p\le19n/2$ and $p\le9n$, respectively.  This covers every residue,
including all integer endpoints.

\paragraph{The local kernel and the cancellation.}
Reduce the two monic polynomials modulo $p$.  The multiplicity inequality
implies that
\[
 \mathcal Q(X)=(X^p-X)\frac{\overline{\mathcal N}(X)}
                               {\overline{\mathcal D}(X)}
             \in\mathbb F_p[X],\qquad \deg\mathcal Q=p-2.
\]
At a residue outside (S), $\mathcal Q(k)=0$.  At a residue inside (S),
substitution $t=k+pV$, cancellation of the factors not divisible by $p$,
and $(X^p-X)'=-1$ give
\[
 pF_n(k+pV)\equiv-\mathcal Q(k)\,
       \frac{V^3}{(V+1)^2(V+2)^2}\pmod p.               \tag{LK}
\]
Outside (S) the left side reduces to zero.  This is an identity of local
rational functions.  All distinct high-digit roots differ by at most
four, and hence by a $p$-adic unit: the interval is empty for $n=1$, and
otherwise $p\ge17$.  It is consequently legitimate to extract the
partial-fraction coefficients from (LK).  It also shows at every residue
that $a_j/p$ and $b_j$ are $p$-integral, and that both reduce to zero away
from the support (S).

On that support, the two nodes in the partial fractions of $F_n$ are
$j=p-k$ and $j=2p-k$.  Their harmonic indices $j+2n$ lie respectively
below $p$ and between $p$ and $2p$.  Thus, after multiplying an order-$r$
harmonic number by $p^r$, their residues modulo $p$ are $0$ and $1$.
The fixed high-digit kernel is
\[
 \frac{V^3}{(V+1)^2(V+2)^2}
 =-\frac1{(V+1)^2}+\frac5{V+1}
   -\frac8{(V+2)^2}-\frac4{V+2}.
\]
Its harmonic readout constants are therefore $-8-4/2=-10$ for $P_3$ and
$-8-4/3=-28/3$ for $P_4$.  Summing (LK) over all residues yields
\[
 p^2P_{3,n}\equiv10\sum_{k\in\mathbb F_p}\mathcal Q(k)=0\pmod p,
 \qquad
 p^3P_{4,n}\equiv\frac{28}{3}
       \sum_{k\in\mathbb F_p}\mathcal Q(k)=0\pmod p.
\]
The last equalities use $\deg\mathcal Q=p-2<p-1$.
Also $Q_n=\sum a_j\in p\mathbb Z_{(p)}$.  This proves (BS).
For the normalization $C_n=(15n)!^2/(10n)!^3$, one has
$v_p(C_n)=-1$ on $(8n,10n)$ and $v_p(C_n)=2$ on $(10n,11n)$.
Accordingly the common exponent in
$6C_nd_{19n}^r(Q_n,P_{r,n})$, $r=3,4$, is at least $1$ and $4$ on these
two intervals, respectively: exactly one more than the baseline local
coefficient bound.

\section{Fourteen first-order shells and all their periodic lifts}
Let $\mathcal S\subset(0,1)$ be the union of the fourteen open intervals
in the first column of the table below.  For every prime $p>29$ with
$p^2>30n$ and $\{n/p\}\in\mathcal S$, one has
\[
 v_p(Q_n)\ge1,\qquad v_p(P_{3,n})\ge-1,\qquad
 v_p(P_{4,n})\ge-2.                                    \tag{PS}
\]
Thus the extra prime divisor may include every such prime, not only those
in the original shell $8n<p<11n$.  These are first-order divisibility
statements.  No unproved second-order enhancement is included.

\paragraph{The finite root-count certificate.}
Put $n/p=t+f$ with $t\in\mathbb Z_{\ge0}$ and $0<f<1$, and write
$y=k/p$ for a residue $0\le k<p$.  At $t=0$, the three numerator high-root
intervals have integer endpoints
\[
 [\lfloor uf-y\rfloor+1,\ \lfloor(u+10)f-y\rfloor],\quad u=0,1,2,
\]
and the two denominator intervals have endpoints
\[
 [\lceil-17f-y\rceil,\ \lfloor-f-y\rfloor],\qquad
 [\lceil-16f-y\rceil,\ \lfloor-2f-y\rfloor].
\]
Let $h$ be the total numerator count minus the total denominator count,
and put $\epsilon_0=\lfloor2f-y\rfloor$.
For every $f\in\mathcal S$ away from Farey walls of order $29$ (in particular, every admissible $f=\{n/p\}$), the following two assertions hold:
$h\ge-1$ at every residue; and wherever $h=-1$, subtracting $\epsilon_0$
from every one of the ten endpoints yields the single row of the table
corresponding to $f$.  Notation $[L,U]$ includes every integer from $L$ to
$U$ with multiplicity; overlapping intervals retain their multiplicities.
\[
\begin{array}{c|ccc|cc}
 f\text{ interval}&N_0&N_1&N_2&D_0&D_1\\\hline
 (1/11,1/8)&[1,1]&[1,1]&[1,1]&[-1,0]&[-1,0]\\
 (3/13,4/17)&[0,1]&[0,2]&[1,2]&[-4,-1]&[-4,-1]\\
 (3/11,5/18)&[0,2]&[0,2]&[1,2]&[-5,-1]&[-4,-1]\\
 (4/13,5/16)&[0,2]&[1,3]&[1,3]&[-5,-1]&[-5,-1]\\
 (4/11,7/19)&[0,3]&[1,3]&[1,4]&[-6,-1]&[-6,-1]\\
 (5/13,7/18)&[1,3]&[1,4]&[1,4]&[-6,-1]&[-6,-1]\\
 (5/12,8/19)&[0,3]&[1,4]&[1,4]&[-7,-1]&[-7,-2]\\
 (5/11,8/17)&[0,3]&[0,4]&[1,4]&[-8,-2]&[-8,-2]\\
 (6/11,9/16)&[0,4]&[0,5]&[1,5]&[-10,-2]&[-9,-2]\\
 (7/12,10/17)&[-1,4]&[0,5]&[1,5]&[-11,-2]&[-10,-3]\\
 (7/11,11/17)&[-1,5]&[0,5]&[1,6]&[-12,-2]&[-11,-3]\\
 (7/10,12/17)&[-1,5]&[0,6]&[1,7]&[-13,-2]&[-12,-3]\\
 (8/11,11/15)&[-1,5]&[0,6]&[1,7]&[-13,-3]&[-13,-3]\\
 (10/11,16/17)&[-1,7]&[0,8]&[1,9]&[-17,-3]&[-16,-4]
\end{array}
\]
This finite table can be verified entirely with rational arithmetic.
All changes of an endpoint occur on a line
$y=mf+c$, with $m\in\{0,1,2,10,11,12,-1,-2,-16,-17\}$ and integer $c$.
The crossing abscissas in $(0,1)$ therefore have reduced denominator at
most $29$.  Partition the $f$-axis by the Farey sequence of order $29$;
there are $270$ open cells.  At the midpoint of each, test every
$y$-boundary and every intervening open interval, including $y=0$ and
excluding $y=1$.  The root data are constant on each tested cell, so this
finite verification proves the table, rather than testing selected
values of $n$.  The supplied certificate checks all $5400$ such cells
and boundaries; $47$ of the Farey cells form the fourteen displayed
intervals.  A missing Farey boundary cannot occur for $f=\{n/p\}$ when
$p>29$; its reduced denominator would have to be $p$.

\paragraph{Why the same table proves every lift.}
Changing $f$ to $t+f$ adds $30t$ numerator roots and $30t$ denominator
roots at every residue.  Hence $h$ and its minimal support are unchanged.
Also $\epsilon=\lfloor2(t+f)-y\rfloor=2t+\epsilon_0$.
For a table entry $N_u=[L_u,U_u]$, the high-root interval after subtracting
$\epsilon$ is
\[
 [(u-2)t+L_u,\ (u+8)t+U_u].
\]
For $D_0=[L_3,U_3]$ and $D_1=[L_4,U_4]$ it is respectively
\[
 [-19t+L_3,-3t+U_3],\qquad[-18t+L_4,-4t+U_4].
\]
These five intervals depend only on $t$ and the table row, and are
independent of the residue $k$ on the entire minimal support.  Let
$U_t(V)$ be the monic rational function having these numerator and
denominator roots.  It has degree $-1$.

As in the broad-shell proof,
$\mathcal Q=(X^p-X)\overline{\mathcal N}/\overline{\mathcal D}$ is a
polynomial of degree $p-2$.  At a minimal-support residue the local
factorization is
\[
 pF_n\bigl(k+p(V+\epsilon)\bigr)
       \equiv-\mathcal Q(k)U_t(V)\pmod p;
\]
outside that support its reduction is zero.  Distinct high roots differ
by at most $\lfloor29n/p\rfloor<p$, so all local partial-fraction
operations have $p$-integral coefficients.  A denominator root $l$ of
$U_t$ has harmonic index $-l$ after extracting the multiples of $p$ from
$H_{j+2n}^{(r)}$.  These indices are at most $19n/p<p$, so their harmonic
numbers are $p$-integral too.  The resulting leading local readout is
therefore a constant depending on $t$ and the table row, multiplied by
$\mathcal Q(k)$.  The sum over $k$ vanishes since
$\deg\mathcal Q<p-1$.  This proves (PS) for every periodic lift.

The normalization valuation
$c(f)=2\lfloor15f\rfloor-3\lfloor10f\rfloor$ is unchanged by the lift.
On the table intervals the baseline profile is exactly $c(f)+1$ for both
readouts, including every lifted interval; (PS) raises it to $c(f)+2$.
The same rational certificate checks these equalities against the
baseline chamber table.

\paragraph{Arithmetic density.}
If
\[
 H_n^*=\prod_{\substack{p>29,\ p^2>30n\\\{n/p\}\in\mathcal S}}p,
\]
then the prime number theorem gives
\[
 \lim_{n\to\infty}\frac{\log H_n^*}{n}
 =\sum_{(a,b)\subset\mathcal S}\sum_{t=0}^{\infty}
      \left(\frac1{t+a}-\frac1{t+b}\right).
\]
The discarded primes have logarithmic total $o(n)$.  All terms are
positive.  At $t=0$ the exact sum is $195197/55440$; summing only
$t=0,\ldots,999$ gives the exact rational lower bound recorded by the
certificate, and in particular
\[
 \liminf_{n\to\infty}\frac{\log H_n^*}{n}
       >3.6747281958.
\]
Thus the previously certified baseline density $26.2366642995$ improves
to a certified density greater than $29.9113924953$.


\section{The divisor density and the reduced denominator}
Define $G_n$ by taking the baseline exponents $g(n/p)$ from (G), and the
additional certified shell exponents, at the eligible primes (E).
The preceding arguments prove
\[
 \frac{6C_nd_{19n}^3}{G_n}(Q_n,P_{3,n})\in\Z^2,
 \qquad
 \frac{6C_nd_{19n}^4}{G_n}(Q_n,P_{3,n},P_{4,n})\in\Z^3.
 \tag{I}
\]
For $B\ne0$, form the integer numerator $AP_{3,n}+BP_{4,n}$ in the second
normalization. For $B=0$, use the first one. Consequently the reduced
denominator actually divides the absolute value of the corresponding
integer multiple of $Q_n$. This is why the argument estimates the
reduced denominator, not the height of an unreduced matrix entry.

By the prime number theorem, the baseline logarithmic density is
\[
 \gamma_0=\int_0^\infty\frac{g(x)}{x^2}\,dx
 =6+\sum_{(a,b,m)}m\sum_{t=0}^\infty
       \left(\frac1{t+a}-\frac1{t+b}\right),
\]
where $(a,b,m)$ ranges over the nonzero intervals of $\psi$ in the
appendix. The initial $6$ is the integral of $h_0(x)/x^2$.
Summing the positive rational terms for $0\le t<1000$ gives, exactly,
\[
 \gamma_0>26.2366642995124>26.2366642995.
\]
The omitted tail is between $0$ and $3/1000$. This is a rational
inequality checked by \texttt{certify\_density.py}; it is not a numerical
prime-density fit. Primes dividing $n$ contribute $O(\log n)$, and primes
below $\sqrt{30n}$ contribute $o(n)$, so the eligibility restrictions do
not change this density.
The additional shell density is $>3.6747281958$, hence
\[
 \liminf\frac{\log G_n}{n}>29.91139.
\]
Since $\log d_{19n}=19n+o(n)$ and
$\log C_n=30n\log(3/2)+o(n)$, (I) gives
\[
 \limsup\frac{\log q_N(A,0)}n
       <122.451536565866-29.91139,
\]
\[
 \limsup\frac{\log q_N(A,B)}n
       <141.451536565866-29.91139\qquad(B\ne0).
\]
Combining these upper bounds with the accepted error exponent
$E>112.807539247454$ proves the theorem by $\delta\ge E/H-1$.
The fixed factors $6,A,B$ contribute no exponential cost. No
irrationality-measure upper bound is asserted: the present conclusion
concerns this approximation sequence and its reduced denominators.

\section{What has and has not been recovered}
At the previously checked depth $N=500$, the actual reduced approximants
had $\delta_N(1,1)=0.0168532791698\ldots$ and
$\delta_N(1,0)=0.2277787154690\ldots$. The proved lower bounds above are
weaker. They explain a positive joint delta and most of the stronger
pure-$\zeta(3)$ arithmetic, but do not identify the exact asymptotic gcd.
In particular, further cancellation suggested by a finite prime sample
is kept out of $G_n$ until it has a uniform proof.

The executable check reconstructs native products, verifies the integer
forms (I), and verifies that the reduced denominator for each tested
coefficient pair divides the certified integer denominator. The checked
values are $n=1,2,3,10,49,99,249,499$, with pairs
$(1,0),(0,1),(1,1),(1,-1),(2,3),(-7,5)$. These checks supplement the
uniform argument and do not replace it. The carry partition is separately
verified using exact rational arithmetic, and the shell proof has an
independent audit.

\paragraph{Related arithmetic methods.}
The elementary denominator argument has a direct precedent in
W.~Zudilin, \emph{Well-poised hypergeometric service for Diophantine problems
of zeta values}, Journal de Th\'eorie des Nombres de Bordeaux
15 (2003), 593--626, especially equations (14)--(16):
\url{https://www.numdam.org/item/JTNB_2003__15_2_593_0.pdf}.
There the basic ingredients are integer-valued binomial polynomials,
their derivatives after lcm multiplication, and factorial-normalized
reciprocal blocks with integer residues. These ingredients justify the
baseline arithmetic above; they do not automatically supply the stronger
common divisor needed for this particular trajectory.

C.~Krattenthaler and T.~Rivoal's \emph{Hyperg\'eom\'etrie et fonction z\^eta de
Riemann}, Memoirs of the AMS 186 (2007), no.~875,
\url{https://arxiv.org/abs/math/0311114}, develops stronger denominator
results using hypergeometric identities. Its general theme is relevant,
but no theorem from that paper is being asserted to apply to this exact
parameter set without checking its hypotheses.

For current qualitative context, Li Lai's preprint
\emph{The linear independence of $1$, $\zeta(3)$, and $\zeta(4)$},
submitted 5 October 2026,
\url{https://arxiv.org/abs/2610.06838}, states precisely that linear
independence result. Lemma 2.2 uses double- and simple-pole coefficients
and harmonic terms of orders two through four, closely matching the
readout structure here. Its method uses determinants and residue-class
arithmetic. We have checked the statement and relevant formulas, not
independently audited its full proof; this work makes no novelty claim
for the qualitative irrationality of every fixed mixed combination.

\appendix
\section{The exact periodic carry table}
The function $\psi$ is zero at rational endpoints (these primes are
excluded by (E)); on each open interval it has the value below.
For every other $x$, equation (G) determines the full baseline exponent.
\begin{center}\small
\begin{tabular}{ccc@{\qquad}ccc}\toprule
$a$&$b$&$\psi$&$a$&$b$&$\psi$\\\midrule

$0$&$1/15$&$0$&$3/5$&$2/3$&$1$\\[0pt]

$1/15$&$1/13$&$1$&$2/3$&$11/16$&$3$\\[0pt]

$1/13$&$1/12$&$2$&$11/16$&$9/13$&$2$\\[0pt]

$1/12$&$1/10$&$3$&$9/13$&$7/10$&$3$\\[0pt]

$1/10$&$2/15$&$0$&$7/10$&$11/15$&$0$\\[0pt]

$2/15$&$2/13$&$1$&$11/15$&$14/19$&$2$\\[0pt]

$2/13$&$3/17$&$2$&$14/19$&$3/4$&$1$\\[0pt]

$3/17$&$2/11$&$1$&$3/4$&$4/5$&$2$\\[0pt]

$2/11$&$1/5$&$2$&$4/5$&$13/16$&$1$\\[0pt]

$1/5$&$4/15$&$1$&$13/16$&$9/11$&$0$\\[0pt]

$4/15$&$5/17$&$3$&$9/11$&$16/19$&$1$\\[0pt]

$5/17$&$3/10$&$2$&$16/19$&$13/15$&$0$\\[0pt]

$3/10$&$1/3$&$0$&$13/15$&$7/8$&$2$\\[0pt]

$1/3$&$2/5$&$2$&$7/8$&$15/17$&$3$\\[0pt]

$2/5$&$7/15$&$1$&$15/17$&$9/10$&$2$\\[0pt]

$7/15$&$9/19$&$3$&$9/10$&$14/15$&$0$\\[0pt]

$9/19$&$1/2$&$2$&$14/15$&$17/18$&$2$\\[0pt]

$1/2$&$8/15$&$0$&$17/18$&$18/19$&$1$\\[0pt]

$8/15$&$7/13$&$1$&$18/19$&$1$&$0$\\[0pt]

$7/13$&$3/5$&$2$&&&\\[0pt]

\bottomrule\end{tabular}\end{center}
\section{Reproducible certificates}
The bundle preserves the \texttt{work/} and \texttt{results/} layout.
Run the following from its root, using Python with \texttt{python-flint}
and \texttt{sympy} installed:
\begin{verbatim}
python work/meijerg5255_arithmetic/profiles/exact_chambers.py
python work/meijerg5255_arithmetic/profiles/certify_density.py
python work/meijerg5255_arithmetic/arithmetic/all_shell_cells.py
python work/meijerg5255_arithmetic/verify_certificate.py
python work/meijerg5255_arithmetic/rates/certify_spectral_constants.py
\end{verbatim}
The chamber and density checks themselves use only Python's standard
library and exact fractions. Native multiplication uses rational Flint
matrices. The scripts do not require the archive notebooks.
\end{document}
